%!TEX root = ../Thesis.tex
\chapter{Conclusion and Future Work}
 Throughout my PhD and this thesis I have attempted to answer both methodological and more applied questions related to uncertainty quantification and calibration of machine learning models. Specifically, I have attempted to answer (in part) the question of \textit{how} to get calibrated models with good predictive uncertainty and \textit{why} it matters for downstream applications. I will here summarize the research questions answered and contributions made in this thesis. \\

 \noindent \textbf{Research Question 1:} \textit{Does optimising classification and feature extraction layers at the same time lead to poorly calibrated neural networks and does imposing structure on the feature embedding space further improve calibration?}\\

\noindent Unlike the other works presented in this thesis, this research question was completely post-rationalized. Nevertheless, the project which started with another goal, resulted in our ICML paper "Decoupling Feature Extraction and Classification Layers for Calibrated Neural Networks" presented in Chaper \ref{chap:ICML} and is perhaps in fact the paper I present in the thesis with the strongest empirical results. In this work we showed that training classifiers end-to-end in part leads to high overconfidence of the trained model and proposed Two-Stage Training (TST) as a remedy to this. We showed that for over-parametrized neural networks, TST was highly effective in increasing predictive uncertainty and model calibration on in-distribution data, at no cost of accuracy and very low training cost. We also proposed another Two-Stage Training method called Variational Two-Stage Training (V-TST) in which one of the hidden layer outputs are treated as a latent variable allowing us to impose structure on the feature embedding space. We showed that training with ELBO in the Stage 2 using a Gaussian prior on the latent variable further improved calibration in some cases. \\

\noindent At the time of publication, this work presented limited theoretical explanations for why these methods work. As previously mentioned, since publication, there are three particular avenues of investigation that could be interesting to pursue to why these these methods work. For TST, it would be interesting to explore the connection between early-stopping in Stage 2 and optimality of weight decay on the last layer. It is possible that one could achieve similar results to TST when training end-to-end by imposing the perfect amount of weight decay on the final layer of the network. For V-TST, it could be interesting to cast the training method as a self-learnt data augmentation method. One could backtrace the sampled $\mathbf{z}$s during training to the input space $\mathbf{x}$ to visualize what the sampled latents correspond to in input space, and perhaps investigate if different variational distributions 
$q(\mathbf{z})$ would lead to different styles of input space corruptions. Finally, applying these methods to regression problems and seeing if they afford similar improvements in calibration would also be interesting. \\

\noindent A final discussion point around these methods is whether or not the way we currently train neural networks from scratch and end-to-end such as in most classification and regression tasks is perhaps a large source of poor uncertainty quantification. If one of the core issues for poor neural network calibration on in distribution data in supervised tasks is the joint learning of features and classification/regression layers, one could perhaps use our methods as motivation for always performing self-supervised learning to learn features, and then learning the supervised task afterwards.  \\

  \noindent \textbf{Research Question 2:} \textit{Does local posterior structure improve Deep Ensembles?} \\

\noindent For Research Question 2, we set out to systematically investigate if including local posterior structure in deep ensembles (DEs) would improve uncertainty quantification and calibration of regular DEs resulting in the paper presented in Chapter \ref{chap:AISTATS}. We did this by applying approximate inference methods to individual MAP estimated models and ensembling them, creating ensembles of BNNs (DE-BNNs). We did this across a variety of data modalities, network architectures and ensemble sizes, and firstly found that generally increasing ensemble sizes improved uncertainty quantification. However, to our surprise, we found that even though DE-BNNs outperformed DEs on both in-distribution and out-of distribution metrics for small ensemble sizes, this performance improvement disappeared when the ensembles grew large enough. In fact, for large ensembles even though DE-BNNs consistently were marginally better on out-of-distribution metrics, DEs on the contrary were consistently better on in-distribution metrics - a finding that shines a new light on DE-BNNs. Based on these results we provided advice to practitioners and briefly discussed new avenues of research for multimodal approximate inference. \\

\noindent I will admit that there are a few leftover questions that I wish we had been able to answer more concretely in this work. We had a range of ideas and intuitions for \textit{why} DE-BNNs fail to outperform DEs when they grow large enough, but never reached a conclusive answer, neither theoretically nor empirically - this is perhaps my biggest regret of this work. To me, the results we observed are extremely hard to accept from an intuitive perspective and understanding of marginalization. Essentially: if individual members of a Gaussian mixture model better approximate the true posterior distribution at each of these modes, how can a mixture model of point estimates (placed at the same modes) perform better (when marginalized over) than the Gaussian mixture model? From my understanding of Bayesian probability theory this seems highly counterintuitive (and possibly cannot be true) which leads me to question if perhaps the local Gaussian approximations themselves are overly crude. A natural counter question to such a statement would be how we then can observe improvements when we use Gaussian approximations on MAP models. To this, the only answer I have is that the Gaussian approximations perhaps regularise the predictions of our networks without necessarily approximating the true posterior. These answers and questions however are not satisfactory and quite hand-wavy, but with more time I would have liked to investigate this closer and perhaps discuss this with more of the Bayesian deep learning community.\\

\noindent \textbf{Research Question 3:} \textit{To what degree does surrogate model uncertainty matter in Bayesian Optimisation?}\\

\noindent For Research Question 3, we wanted to empirically investigate the widespread consensus that good calibration of the surrogate model used in Bayesian Optimisation (BO) is important for good performance. The answering of this research question resulted in a UAI paper which I presented in Chapter \ref{chap:UAI}. In this work we showed on a large set of optimisation problems, that across different surrogate model types there is a correlation between model calibration and BO performance, i.e. better calibrated surrogate model types generally perform better in BO. Interestingly, this correlation disappeared when controlling for the surrogate model type, an indication that practitioners should mainly be concerned with picking surrogate model types that generally are well calibrated, but that later optimising for good calibration may not be necessary. We also showed that Deep Ensembles can work surprisingly well as surrogate models. Finally, we showed both theoretically and empirically that using previously proposed re-calibration methods as proposed by \citet{recalibration_BO} during BO may in fact lead to \textit{worse} BO performance. \\

\noindent Since we published this work, \cite{recalibration_BO} have updated their proposed methods for re-calibrating and may now in fact be more useful during BO. In addition, others have performed large works on BNNs as surrogate models \citep{li2024a}, but reach different conclusions to ours - they do not find that deep ensembles perform particularly well as surrogate models which could be both due to architecture choices or BO setup. It would be interesting to dive into exactly how one designs good deep ensemble surrogates for BO, and investigate when they perform well. \\


 \noindent \textbf{Research Question 4:} \textit{How do Variational Bayesian Last Layer (VBLL) neural networks perform as surrogates in Bayesian Optimisation and how can we improve fitting time of these surrogates in this setting?} \\

 \noindent For Research Question 4 we were interested in applying Variational Bayesian Last Layers (VBLLs) on neural networks for surrogate model usage in Bayesian Optimisation. We applied these models to several optimisation problems, both synthetic and non-synthetic structured problems and surprisingly found that they regularly performed on par with GPs on synthetic problems, and outperformed them on problems with structure. Moreover, we showed that Thompson Sampling lends itself especially well to Bayesian Last Layer models often affording strong BO performance. Due to high surrogate fit times of VBLL models, we also proposed a continual learning scheme which interleaves classical stochastic gradient descent optimisation with recursive updating typically associated with Bayesian linear regression, resulting in significantly lower fit times although at a small cost of BO performance. \\

 \noindent From our ICLR work, there are two particular avenues of further work that I find highly interesting. One is something we discussed among the authors, which is designing more principled ways of performing continual learning. In the current continual learning scheme, we set regular intervals of when to perform classical neural network fitting and when to perform Bayesian conditioning. An adaptive scheme for when to perform the two steps could both alleviate fit times and increase performance if designed well. Additionally, we also experimented briefly with using MAP estimated neural networks as surrogate functions, essentially treating them as Thompson Samples which are randomized due to the random initialization of the network at every iteration of BO, and found that these surrogates actually performed very well - another interesting avenue to investigate. \\

\noindent \textbf{Research Question 5:} \textit{How can we improve e-learning by utilizing probabilistic and Bayesian modelling?}\\

\noindent Research Question 5 stemmed from my supervisor Michael's collaboration with the WriteReader startup which has an e-learning platform. The NLDL paper "Neural machine translation for
automated feedback on children’s
early-stage writing" and the pre-print "Calibrated Transformer Models for
Supporting Children’s Writing in
Low-Resource Language" both aimed at answering this research question. In the NLDL paper, we presented a robust likelihood for stabilizing Transformer fine-tuning in English, when the used training data is highly noisy, resulting in better calibrated models. Moreover, we showed that ensembling these models further improved calibration. We showed how the likelihoods of predicted sequences using these models could also be used to reject particularly poor translations which could be used in human-in-the-loop systems. In the pre-print we trained Transformers from scratch in Danish, and showed how different tokenizers influenced calibration metrics. Moreover, we showed that BPE tokenizers generally provided better calibrated models, which again could be useful in human-in-the-loop systems. Finally we showed how these models could be useful in automatic proficiency scoring pipelines. \\

\noindent Both of these works aim at creating models for usage in an e-learning platform, and have so far been tested in sand-box environments which of course is the first natural testing environment. Having said that, given the level of performance in these sand-box environments, the ideal progression for these projects would be to go to human studies, i.e. having test subjects (teachers) evaluating the translations (some of which has already been done for other models in WriteReader). Additionally, it would also be of interest to use some of the newly available pre-trained models in Danish such as those presented on \url{https://www.foundationmodels.dk/}. This was previously not possible, as there were no large-scale pre-trained decoder models in Danish, but Munin 7B Alpha is now available for Danish. Finally, there is also a natural opportunity for applying additional BNN methods on the trained Transformers. \\

\noindent To conclude, in Papers A and B, I investigated methods for calibrated machine learning models: in particular, methods for improving calibration and uncertainty quantification of neural networks. In Papers C and D, I investigated the impacts of surrogate models in BO: in particular, how model calibration impacted BO performance, and applied a new class of Bayesian Neural Networks as surrogate models along with methodological proposals for improving fitting time of such models. And finally, in Papers E and F, I investigated using Transformers for a sequence-to-sequence task and showed how the inherent uncertainty quantification and calibration of these Transformers would impact and could be used in a downstream application. As such, I have presented and concluded on the contributions made in this thesis titled \textit{Calibrated Machine Learning Models: How To Get Them and Why They Matter}. I thank you for your time.